\subsubsection{Dirichlet 变量口径（与 Abel 单位圆盘完全等价）}\label{app:real-input-40-dirichlet}
为对接主文第 \ref{sec:zeta-finite-part} 节的"轨道 Dirichlet 级数/Abel 有限部"语言（定义 \ref{def:abel-finite-part}），
可把 primitive 生成函数 $P$ 改写为 $s$-变量的轨道 Dirichlet 级数。定义
\[
\mathcal{P}(s):=\sum_{n\ge 1}p_n\,\lambda^{-ns}=P(\lambda^{-s}),
\qquad
\mathcal{Z}(s):=\zeta_M(\lambda^{-s}).
\]
注意到 Abel 归一化的 $r=\lambda z$ 与 Dirichlet 变量满足
\[
z=\lambda^{-s},\qquad r=\lambda^{1-s}\in(0,1)\quad(s>1),
\]
因此 $\mathcal{Z}(s)=\widetilde{\zeta}_M(\lambda^{1-s})$，而 $\mathcal{P}(s)$ 亦等价于 $\mathcal{P}(r)$ 在 $r=\lambda^{1-s}$ 的替换。
由于
\[
1-\lambda^{1-s}=1-e^{-(s-1)\log\lambda}\sim (s-1)\log\lambda\qquad(s\to 1^+),
\]
由 Abel 口径的奇性展开
\(
\mathcal{P}(r)=-\log(1-r)+\log\mathfrak{M}+o(1)
\)
得到等价的 Dirichlet 口径：
\[
\mathcal{P}(s)=-\log(s-1)-\log\log\lambda+\log\mathfrak{M}+o(1)\qquad(s\to 1^+).
\]
从而 $s$-平面上的 Hadamard 有限部常数满足
\[
\mathrm{fp}_{s=1}\bigl(\mathcal{P}(s)+\log(s-1)\bigr)=\log\mathfrak{M}-\log\log\lambda.
\]
另一方面，由 $\widetilde{\zeta}_M(r)=\frac{C}{1-r}+O(1)$ 与 $1-r\sim (s-1)\log\lambda$ 可得
\[
\mathcal{Z}(s)\sim \frac{C}{(s-1)\log\lambda}\qquad(s\to 1^+).
\]
这给出 Abel 单位圆盘变量与 Dirichlet 变量两套口径的完全等价闭式接口。

\paragraph{盘内极坐标的角变量（无参口径）}
上述等价把临界奇性完全收束在径向参数 \(r\to 1^{-}\) 上；若需要在同一 Abel 圆盘内同时组织一个\emph{角}变量，可取附录 \ref{app:unit-circle-phase-gate} 的无参面积相位紧化：对任意 \(x\in\RR\) 令
\[
e^{\mathrm{i}\vartheta(x)}=\iota(x)=\frac{1+\mathrm{i}x}{1-\mathrm{i}x}\in\TT,
\qquad
\vartheta(x)=2\arctan(x),
\]
则盘内点可统一写成 \(w=r\,e^{\mathrm{i}\vartheta(x)}\)（见 \eqref{eq:unit-circle-phase-disk-polar}）。当 \(r=1\) 时退化为 unitary slice 的无参参数化 \(w=\iota(x)\)；该补充不改变 Dirichlet--Abel 的径向奇性分析，仅提供一个与 unitary-slice 叙述对齐的角变量来源。

\paragraph{Dirichlet--Witt 伸缩算子（严格交换的桥包）}
由 Euler 乘积/primitive 分解恒等式
\(
\log\zeta(z)=\sum_{m\ge 1}\frac{1}{m}P(z^m)
\)
作代换 $z=\lambda^{-s}$，注意到 $z^m=\lambda^{-ms}$，得到
\[
\log\mathcal{Z}(s)=\sum_{m\ge 1}\frac{1}{m}\,\mathcal{P}(ms).
\]
因此在 Dirichlet 变量上引入 Witt--dilation（伸缩）算子
\[
(Wf)(s):=\sum_{m\ge 1}\frac{1}{m}\,f(ms),
\]
则有严格恒等式
\[
\boxed{\ \log\mathcal{Z}(s)=(W\mathcal{P})(s)\ }.
\]
对该等式作 Möbius 反演（即 $W^{-1}$）得到 primitive 的 Dirichlet 级数可由 $\mathcal{Z}$ 直接剥离：
\[
\boxed{\ \mathcal{P}(s)=\sum_{m\ge 1}\frac{\mu(m)}{m}\,\log\mathcal{Z}(ms)\ }.
\]
这两条等式把“谱/迹侧”的 $\mathcal{Z}(s)=\det(I-\lambda^{-s}M)^{-1}$ 与 “prime-like 侧”的 $\mathcal{P}(s)$ 放在同一接口下严格对齐；其中 $s\mapsto ms$ 的伸缩正是 Abel 变量的幂 $z\mapsto z^m$ 在 Dirichlet 口径中的等价翻译。

\paragraph{Witt 伸缩的无参 Koenigs 线性化坐标：伸缩 \texorpdfstring{$(s\mapsto ms)$}{(s->ms)} 的平移化}
在 Dirichlet--Abel 接口 \(z=\lambda^{-s}\) 下，伸缩 \(s\mapsto ms\) 与幂替换 \(z\mapsto z^{m}\) 严格等价。
因此可将 Witt 伸缩链沿“对数的对数”坐标线性化为平移半群。

\begin{theorem}[Koenigs 线性化：伸缩 \texorpdfstring{$(s\mapsto ms)$}{(s->ms)} 的平移化]\label{thm:koenigs-linearization-witt-dilation}
\leanverified{paper\_real\_input\_40\_koenigs\_linearization\_witt\_dilation}
\leanverified{paper\_koenigs\_linearization\_witt\_dilation}
在实域 \(s>0\) 上定义无参深度坐标
\begin{equation}\label{eq:witt-koenigs-depth-coordinate}
\boxed{\;
y(s):=\log\!\bigl(s\log\lambda\bigr)
=\log\!\bigl(-\log(\lambda^{-s})\bigr)\; }.
\end{equation}
则对任意整数 \(m\ge 1\) 有严格恒等式
\begin{equation}\label{eq:witt-koenigs-translation}
\boxed{\; y(ms)=y(s)+\log m\; }.
\end{equation}
从而 Witt 伸缩 \(s\mapsto ms\) 在 \(y\)-坐标下被完全对角化为平移半群 \(y\mapsto y+\log m\)，且该线性化不引入外部尺度参数（\(\lambda\) 已由 Dirichlet--Abel 接口固定）。
\end{theorem}
\begin{proof}
由 \eqref{eq:witt-koenigs-depth-coordinate} 得
\[
y(ms)=\log\!\bigl(ms\log\lambda\bigr)=\log\!\bigl(s\log\lambda\bigr)+\log m=y(s)+\log m,
\]
即得 \eqref{eq:witt-koenigs-translation}。证毕。
\end{proof}

\begin{corollary}[Witt 伸缩算子在深度线上的平移卷积表示]\label{cor:witt-depth-line-convolution}
\leanverified{paper\_witt\_depth\_line\_convolution}
令 \(f:(0,\infty)\to\CC\)，并令
\[
s=\frac{e^{y}}{\log\lambda},
\qquad
g(y):=f\!\Bigl(\frac{e^{y}}{\log\lambda}\Bigr).
\]
则 Witt 伸缩算子 \(W\) 满足恒等式
\begin{equation}\label{eq:witt-depth-line-convolution}
\boxed{\;
(Wf)\!\Bigl(\frac{e^{y}}{\log\lambda}\Bigr)
=\sum_{m\ge 1}\frac{1}{m}\,g(y+\log m)\; }.
\end{equation}
等价地，若记平移算子 \(T_{\tau}g(y):=g(y+\tau)\)，则 \(W\) 在深度线 \(y\) 上成为由离散测度
\(
\nu:=\sum_{m\ge 1}\frac{1}{m}\delta_{\log m}
\)
生成的平移叠加算子。
\end{corollary}
\begin{proof}
由定理 \ref{thm:koenigs-linearization-witt-dilation}，
\(
g(y+\log m)=f(\tfrac{e^{y+\log m}}{\log\lambda})=f(ms)
\)。
代入 \( (Wf)(s)=\sum_{m\ge 1}\frac{1}{m}f(ms)\) 即得 \eqref{eq:witt-depth-line-convolution}。证毕。
\end{proof}

\paragraph{正则化深度卷积与素数局部分解：Euler 乘积的算子化}
为获得收敛的符号口径，对 \(\sigma>0\) 定义正则化深度卷积算子
\begin{equation}\label{eq:witt-depth-regularized-operator}
\boxed{\;
\mathcal{K}_{\sigma}
:=\sum_{m\ge 1}\frac{1}{m^{1+\sigma}}\,T_{\log m}\; }.
\end{equation}

\begin{theorem}[算子 Euler 分解：素数作为平移生成元]\label{thm:witt-depth-euler-product}
在任意使 \(T_{\tau}\) 成为等距作用的 Banach 空间上，若 \(\sigma>0\) 使 \(\sum_{p}p^{-(1+\sigma)}<\infty\)，则 \(\mathcal{K}_{\sigma}\) 具有收敛于算子范数的局部分解
\begin{equation}\label{eq:witt-depth-euler-product}
\boxed{\;
\mathcal{K}_{\sigma}
=\prod_{p}\Bigl(\sum_{k\ge 0}p^{-k(1+\sigma)}\,T_{k\log p}\Bigr)
=\prod_{p}\bigl(I-p^{-(1+\sigma)}T_{\log p}\bigr)^{-1}\; }.
\end{equation}
该分解将 Euler 乘积提升为平移半群上的局部可逆分解：素数对应基本平移 \(T_{\log p}\)，而合成尺度对应其几何级数闭包。
\end{theorem}
\begin{proof}
对每个素数 \(p\)，由 \(\|T_{\log p}\|=1\) 得 \(\|p^{-(1+\sigma)}T_{\log p}\|=p^{-(1+\sigma)}<1\)，从而
\[
\bigl(I-p^{-(1+\sigma)}T_{\log p}\bigr)^{-1}
=\sum_{k\ge 0}p^{-k(1+\sigma)}T_{k\log p}
\]
在算子范数意义下收敛。
\par
将有限素数集合上的乘积展开，可写为
\[
\prod_{p\le P}\Bigl(\sum_{k\ge 0}p^{-k(1+\sigma)}T_{k\log p}\Bigr)
=\sum_{m\in\mathcal{M}_P}\frac{1}{m^{1+\sigma}}\,T_{\log m},
\]
其中 \(\mathcal{M}_P\) 为所有素因子不超过 \(P\) 的正整数集合。
由唯一分解 \(m=\prod_{p\le P}p^{k_p}\) 得 \(\log m=\sum_{p\le P}k_p\log p\) 且 \(m^{-(1+\sigma)}=\prod_{p\le P}p^{-k_p(1+\sigma)}\)，故权重与平移参数严格匹配。
令 \(P\uparrow\infty\) 并用 \(\sigma>0\) 导致的绝对收敛交换求和与取极限，即得 \eqref{eq:witt-depth-euler-product}。证毕。
\leanverified{paper\_witt\_depth\_euler\_product}
\end{proof}

\begin{corollary}[有限深度裁剪：有限素数--有限幂级数窗口]\label{cor:witt-depth-finite-window-truncation}
\leanverified{paper\_witt\_depth\_finite\_window\_truncation}
对固定 \(\sigma>0\) 与深度窗口尺度 \(L>0\)，定义有限裁剪算子
\begin{equation}\label{eq:witt-depth-finite-window-trunc-operator}
\boxed{\;
\mathcal{K}_{\sigma}^{(L)}
:=\prod_{p\le e^{L}}
\Bigl(\sum_{0\le k\le \lfloor L/\log p\rfloor}p^{-k(1+\sigma)}\,T_{k\log p}\Bigr)\; }.
\end{equation}
则 \(\mathcal{K}_{\sigma}^{(L)}\) 仅依赖有限个素数与有限个 \(p\)-幂平移层。
并且对任意有界 \(g\)，余项
\(
\mathcal{R}_{\sigma,L}:=(\mathcal{K}_{\sigma}-\mathcal{K}_{\sigma}^{(L)})
\)
仅由（i）素数尾部 \(p>e^{L}\) 与（ii）局部几何级数尾部 \(k>L/\log p\) 产生；其在一致范数下可由
\(
\sum_{p>e^{L}}p^{-(1+\sigma)}
\)
以及
\(
\sum_{k> L/\log p}p^{-k(1+\sigma)}
\)
两类显式尾和给出上界。
\end{corollary}

\begin{remark}[局部 \(p\)-层截断与 Frobenius 缺陷的整数证书]\label{rem:witt-depth-local-p-frobenius}
定理 \ref{thm:witt-depth-euler-product} 将 \(\mathcal{K}_{\sigma}\) 分解为素数局部因子
\[
\bigl(I-p^{-(1+\sigma)}T_{\log p}\bigr)^{-1}
=\sum_{k\ge 0}p^{-k(1+\sigma)}\,T_{k\log p},
\]
其 \(k\)-项对应于 \(p\)-幂尺度 \(p^{k}\) 的深度平移层。
因此将局部因子截断到 \(k\le K\) 等价于仅保留 \(p\)-幂尺度的有限层数据。
\par
另一方面，在系数环版本的 Big-Witt/Möbius 反演中，Adams 幂替换 \(\psi^{p}\) 在多项式环上由 \(u\mapsto u^{p}\) 实现，而层间一致性可由 Frobenius 缺陷
\(
\Delta_{p,k}(u)
\)
（定义 \ref{def:witt-frobenius-defect}）及其整除深度
\(
\Delta_{p,k}(u)\in p^{k}\ZZ[u]
\)
（命题 \ref{prop:witt-frobenius-defect-divisibility}）审计。
当该整除深度在 \(k\le K\) 上成立时，局部 \(p\)-层截断在模 \(p^{K}\) 意义下与 Frobenius 推前兼容；一旦该证书失败，则任何关于端点模类通道（见推论 \ref{cor:endpoint-tristate-compression} 及定义 \ref{def:endpoint-frobenius-integer-certificate}）的进一步谱层解释在结构上均失去实现层支撑。
\end{remark}
